\documentclass[12pt]{article}
\usepackage[utf8]{inputenc}
\usepackage[T1]{fontenc}
\usepackage{hyperref}
\title{Tsou Language Resources for K-12 Curriculum Development}
\author{Compiled by Claude Sonnet 5 (claude-sonnet-5) for Andr\'as Kornai}
\date{}
\begin{document}
\maketitle
\section{Introduction}
Tsou (ISO 639-3 tsu) is an endangered Austronesian language of Taiwan with about 2,100 speakers. The search found only general and academic material. It did not find a named textbook. A Chinese-language follow-up search (for example on the Council of Indigenous Peoples or Ministry of Education) is advisable.
\section{Primary Educational Resources}
No specific textbook or primer was found. Omniglot says revitalization efforts include school teaching and a dictionary, but I could not confirm details: \url{https://www.omniglot.com/writing/tsou.htm}.
\section{Linguistic and Reference Materials}
T'ung-ho Tung's \emph{A Descriptive Study of the Tsou Language} (1991) is catalogued at Academia Sinica: \url{https://ir.sinica.edu.tw/handle/201000000A/6224}. The Austronesian typology database has a language entry: \url{https://lpan.eva.mpg.de/austronesian/language.php?id=824}. A background overview is at \url{https://thelanguagecloset.com/2021/03/06/languages-of-taiwan-tsou/}.
\section{Religious and Community Resources}
None confirmed in search results.
\section{Digital and Audio Resources}
The Endangered Languages Project holds a 1993 Tsou word list: \url{https://endangeredlanguages.com/hi/node/96659}. A blog mentions an online electronic Tsou dictionary, but the search did not give its name or address.
\section{Curriculum Development Context}
Descriptive linguistics for Tsou exists, but this search found no curriculum-ready materials. Official Taiwanese indigenous-language teaching materials probably exist and should be sought in Chinese.
\end{document}